\section{Spearman Correlation using Different Similarity Functions}
\label{section:corr_more}

We present results using TF-IDF-based similarity and Levenshtein similarity when calculating the Spearman correlation. 
Specifically, within each sample of \( n \) proposed answers, 
we calculate Spearman correlation coefficient between the \( n \) similarity scores and the \( n \) preference scores determined by the GPT-4-based evaluator. 
As shown in \Cref{fig:corr_more}, there is indeed a positive correlation between win rate and both TF-IDF similarity and Levenshtein similarity.

\begin{figure}
    \centering
    \begin{subfigure}[b]{0.495\linewidth}
        \includegraphics[height=5cm]{corr_tfidf.pdf}
        % \caption{Caption}
    \end{subfigure}
    % \end{minipage}
    \hfill
    \begin{subfigure}[b]{0.495\linewidth}
    % \begin{minipage}[t]{0.49\linewidth}
        \includegraphics[height=5cm]{corr_edit.pdf}
        % \caption{Caption}
        \label{fig:h2}
    \end{subfigure}
    % \end{minipage}
    \caption{(a) Spearman Correlation using TF-IDF similarity; (b) Spearman Correlation using Levenshtein similarity.}
    \label{fig:corr_more}
\end{figure}

\section{LLM Ranker}

This section introduces the setup of the LLM-Ranker used in this paper. 
The LLM-Ranker is designed to evaluate and rank the best output generated by some LLMs.
\Cref{tab:template_ranking} presents the template for prompting the model during these evaluations. 
We use this LLM-Ranker to pick the best answer among and use AlpacaEval evaluator 
to evaluate the best ranked answer.


\begin{table}[t]
    \centering
    \small
    \caption{Prompt for ranking with LLMs}
    \begin{tabular}{@{}p{1\linewidth}@{}}
    \toprule
    
You are a highly efficient assistant, who evaluates and selects the best large language model (LLMs) based on the quality of their responses to a given instruction. This process will be used to create a leaderboard reflecting the most accurate and human-preferred answers.
\\
I require a leaderboard for various large language models. I'll provide you with prompts given to these models and their corresponding outputs. Your task is to assess these responses, and select the model that produces the best output from a human perspective.
\\
\\
\#\# Instruction
\\
\begin{verbatim}
{
    "instruction": """{instruction}""",
}
\end{verbatim}
\\
\#\# Model Outputs
\\
Here are the unordered outputs from the models. Each output is associated with a specific model, identified by a unique model identifier.

\begin{verbatim}
{
    {
        "model_identifier": "{identifier_1}",
        "output": """{output_1}"""
    },
    {
        "model_identifier": "{identifier_2}",
        "output": """{output_2}"""
    },
    {
        "model_identifier": "{identifier_3}",
        "output": """{output_3}"""
    },
    {
        "model_identifier": "{identifier_4}",
        "output": """{output_4}"""
    },
    {
        "model_identifier": "{identifier_5}",
        "output": """{output_5}"""
    },
    {
        "model_identifier": "{identifier_6}",
        "output": """{output_6}"""
    }
}
\end{verbatim}
\\
\\
\#\# Task
\\
Evaluate the models based on the quality and relevance of their outputs, and select the model that generated the best output. Answer by providing the model identifier of the best model. We will use your output as the name of the best model, so make sure your output only contains one of the following model identifiers and nothing else (no quotes, no spaces, no new lines, ...).
\\
\\
\#\# Best Model Identifier
\\
    \bottomrule
    \end{tabular}
    \label{tab:template_ranking}
\end{table}



\section{Case Study}

\begin{table}[t]
    \centering
    \small
    \caption{Case: Some models produce high quality answers.}
    \begin{tabular}{@{}llp{0.6\linewidth}@{}}
    \toprule
    Role        &  Preference   & Content \\
    \midrule
    User        &      -         & I want to talk about "Smooth",  a song by Rob Thomas. \\
    \midrule
    Qwen1.5-110B-Chat            & 0.35
    &  ... \setulcolor{blue}\ul{It's a perfect blend of rock, pop, and Latin music, showcasing Santana's signature guitar work and Rob Thomas's soulful vocals} ...
       
    \\   
    Qwen1.5-72B-Chat             & 0.00
    &  ...
        
    \\     
    Llama-3-70B-Instruct         & 0.00
    & ... \setulcolor{red}\ul{it's a collaboration between Rob Thomas, the lead vocalist of Matchbox Twenty, and Santana, the legendary guitarist} ...
    
    \\         
    WizardLM-2-8x22B             & 0.99
    & ... \setulcolor{orange}\ul{It spent 12 weeks at the top of the Billboard Hot 100 chart in the United States} ...
         \setulcolor{orange}\ul{capturing the song's sultry and energetic vibe} ...
    
    \\     
    Mixtral-8x22B-Instruct-v0.1  & 0.00
    & ...
    
    \\        
    dbrx-instruc                 & 0.00
    & ...
    
    \\ 
    \midrule
    \begin{tabular}{@{}l@{}}
    Aggregated \\
    (Qwen1.5-110B-Chat)
    \end{tabular}
    & 0.99
    & "Smooth" is actually \setulcolor{red}\ul{a collaboration between Rob Thomas, who is the lead singer of the band Matchbox Twenty, and the renowned Latin rock guitarist Carlos Santana} ... \setulcolor{blue}\ul{This iconic song blends rock, pop, and Latin music elements, showcasing Santana's signature guitar work and Rob Thomas's distinctive vocals} ... \setulcolor{orange}\ul{The song was a colossal success, spending 12 weeks at the top of the Billboard Hot 100} ... The music video, \setulcolor{orange}\ul{set in a vibrant club scene, complements the song's sultry vibe} ...
    
    \\ 
    \bottomrule
    \end{tabular}
    \label{tab:case1}
\end{table}

\begin{table}[t]
    \centering
    \small
    \caption{Case: all proposed responses are not good enough.}
    \begin{tabular}{@{}llp{0.6\linewidth}@{}}
    \toprule
    Role        &  Preference   & Content \\
    \midrule
    User        &      -         & How do you become an author? \\
    \midrule
    Qwen1.5-110B-Chat            & 0.00 
    &  Becoming an author involves several steps and ...
    3. \setulcolor{blue}\ul{**Choose your niche**}:...
    8. \setulcolor{blue}\ul{**Understand the publishing industry**}:...
    10. \setulcolor{blue}\ul{**Self-publishing options**}:...
       
    \\   
    Qwen1.5-72B-Chat             & 0.00
    &  ...
        
    \\     
    Llama-3-70B-Instruct         & 0.16
    & Becoming an author can be a rewarding and fulfilling career...
    1. \setulcolor{red}\ul{**Develop a love for reading and writing**}:...
    10. \setulcolor{red}\ul{**Market and promote your work**}:...
    11. \setulcolor{red}\ul{**Continuously improve**}:...
    
    \\         
    WizardLM-2-8x22B             & 0.03
    & 1. ...
    4. \setulcolor{orange}\ul{Finish your work}:...
    
    \\     
    Mixtral-8x22B-Instruct-v0.1  & 0.00
    & ...
    
    \\        
    dbrx-instruc                 & 0.00
    & ...
    
    \\ 
    \midrule
    \begin{tabular}{@{}l@{}}
    Aggregated \\
    (Qwen1.5-110B-Chat)
    \end{tabular}
    & 0.33
    & Becoming an author involves a dedicated journey...
    1. \setulcolor{red}\ul{**Cultivate a Love for Writing and Reading**}:...
    4. \setulcolor{blue}\ul{**Choose Your Niche**}:...
    5. \setulcolor{orange}\ul{**Finish Your Work**}:...
    10. \setulcolor{blue}\ul{**Self-Publishing**}:...
    11. \setulcolor{red}\ul{**Marketing and Promotion**}:...
    12. \setulcolor{red}\ul{**Continuous Learning and Writing**}:...
    
    \\ 
    \bottomrule
    \end{tabular}
    \label{tab:case2}
\end{table}


We present a case study in this section. 
Due to the length of the responses generated by all models, we will only show selected fragments for brevity. 
To illustrate how the aggregator synthesizes the response, 
we underlined similar expressions between the proposed responses and the aggregated response in different colors.
We omit the content that all proposed responses have mentioned.

\Cref{tab:case1} showcases the responses generated by different proposers.
The aggregated response generated by Qwen1.5-110B-Chat 
reflects a high preference for its own content but also incorporates key points 
from Llama-3-70B-Instruct and WizardLM 8x22B. 
Notably, GPT-4's preference score for WizardLM 8x22B's response is 0.99, 
and the final aggregated answer also achieves a preference score of 0.99.

Meanwhile, \Cref{tab:case2} presents another case
where none of the proposed responses achieve a high GPT-4 preference score. 
Despite this, the aggregator successfully identifies and 
incorporates the strong points from these responses,
achieving a preference score of 0.33.


\section{MATH Task}

Here, we demonstrate that our approach is applicable to reasoning tasks, 
such as those in the MATH dataset \cite{hendrycks2021measuring}.
The results are presented in \Cref{tab:MATH}, where we show that our method consistently enhances accuracy by a significant margin. 
This indicates that our approach is also effective for this type of task.
Notably, our method is complementary to existing reasoning techniques such as Chain of Thought \cite{cot} and Self-consistency \cite{cot-sc}.

\begin{table}[t]
\centering
\caption{
Results on the MATH task. 
We evaluate different aggregators,
with all six models serving as proposers in each MoA layer.
}
\begin{tabular}{@{}lllll@{}}
\toprule
Aggregator                  & Layer 1 & Layer 2 & Layer 3  \\
\midrule
Qwen1.5-72B-Chat            & 0.428              & 0.526   & 0.552    \\
Qwen1.5-110B-Chat           & 0.500              & 0.570   & 0.576    \\
Wizard 8x22b                & 0.544              & 0.574   & 0.580    \\
Mixtral-8x22B-Instruct-v0.1 & 0.282              & 0.534   & 0.556    \\
Llama-3-70B-Instruct        & 0.456              & 0.584   & 0.578    \\
dbrx-instruct               & 0.314              & 0.456   & 0.522    \\
% gpt-4o-2024-05-13           & 0.722              & 0.708   & 0.654   \\
\bottomrule
\end{tabular}
\label{tab:MATH}
\end{table}
